%! @@TITLE@@ — Unit @@UNITINT@@, Lesson @@LESSONID@@ — Homework (Answer Key)
% Mirrors homework/main.tex. Each \writelines{n} is answered with EXACTLY n \ansline's, and every
% work block is authored byte-identically, so the two files come out the same number of pages.
\documentclass[10pt]{article}
\usepackage{@@PREFIX@@-article}
\usepackage{@@PREFIX@@-key}   % pulls in -boxes; do NOT also load -boxes
% Coordinate grid: {xmin}{xmax}{ymin}{ymax}
\newcommand{\lgrid}[4]{%
  \draw[step=1, linegray!40, very thin] (#1,#3) grid (#2,#4);
  \draw[->, semithick, charcoal] (#1,0)--(#2,0) node[below right, font=\scriptsize]{$x$};
  \draw[->, semithick, charcoal] (0,#3)--(0,#4) node[left, font=\scriptsize]{$y$};}

\begin{document}

\pageheader{Unit @@UNITINT@@, Lesson @@LESSONID@@}{Homework --- Answer Key}
% NAMESTRIP: no \namedateperiod here — mirrors the blank.

\vspace{0.10in}

% TODO: copy homework/main.tex verbatim from here down, then fill every blank.
%   \blank{2.0cm}   ->  \ans{answer}
%   \writelines{2}  ->  two \ansline{...}
%   \begin{work} ... \end{work}  ->  UNCHANGED, byte for byte
% For the multiple-choice item, keep all four options and wrap the CORRECT one in \ans{...}, then
% use the answer lines below it to say which is right and why one distractor is wrong.
\begin{notesbox}{Practice}
Show your work. TODO --- identical to the blank.
\end{notesbox}

\end{document}
